\begin{textsample}{NEG Dim 1 – 2011\_c6 – Score -4.00 – 2011\_c6/enade\_2011\_tec\_man\_ind\_da8c9f098c.txt}  \label{ex:f1_neg_227}
CAPÍTULO 3 ANÁLISE TÉCNICA DA PROVA Este capítulo tem por objetivo apresentar o desempenho dos estudantes concluintes de Tecnologia em Manutenção Industrial no ENADE/2011.
Para isso, foram calculadas as estatísticas básicas da prova em seu todo, bem como as estatísticas dos componentes relacionadas à Formação Geral, ao de Conhecimento Específico da Área e das questões discursivas isoladamente.
Nas tabelas, são apresentados o tamanho da população inscrita e de presentes, e as seguintes estatísticas das notas⁸ : média do desempenho na prova, erro padrão da média, desvio padrão, nota mínima, mediana e nota máxima.
As estatísticas apresentadas neste capítulo contemplam o total de estudantes concluintes da área de Tecnologia em Manutenção Industrial em 2011 do Brasil e, separadamente, por Grande Região.
Foram calculadas tendo -se em vista as seguintes agregações : ( a ) as Grandes Regiões e o país como um todo ; ( b ) a Categoria Administrativa ; e ( c ) a Organização Acadêmica.
Em relação aos gráficos de distribuição de notas, o intervalo considerado foi de 10 unidades, aberto à esquerda e fechado à direita, com exceção do primeiro intervalo, [ 0 ; 10 ], fechado em ambos os extremos.
Para os gráficos de distribuição das notas das questões discursivas, foram consideradas mais duas categorias : questão em branco e nota zero. 3.1 ESTATÍSTICAS BÁSICAS DA PROVA 3.1.1 Estatísticas Básicas Gerais A Tabela 3.1 apresenta as estatísticas básicas da prova por grande Região.
A população total de inscritos foi de 498.
Destes, 408 \textbf{estiveram} presentes, sendo 18,1 \% o índice de não comparecimento.
A Região de maior abstenção foi a Centro-Oeste ( 28,8 \% ) e a de menor abstenção foi a Norte ( 3,1 \% ).
A média das notas da prova como um todo ( nas seções seguintes serão analisados os componentes de Formação Geral e de Conhecimento Específico ) foi 42,8, sendo que os alunos da região Norte obtiveram a média mais baixa ( 39,6 ), e os da região Nordeste obtiveram, a média mais alta ( 46,3 ).
As demais médias foram : 42,0 na região Sudeste ; 43,6 na região Sul ; e 42,1 na região Centro-Oeste.
O desvio padrão para o Brasil como um todo foi 12,2, sendo o maior desvio padrão encontrado na região Centro-Oeste ( 12,6 ) e o menor ⁸ Essas estatísticas e outras \textbf{estão} definidas no Capítulo 1. na região Norte ( 10,8 ), indicando uma dispersão um pouco menor das notas desta última região.
A região que obteve a maior nota máxima foi a Sul ( 84,6 ), ao passo que a região que atingiu a menor nota máxima foi a Norte ( 65,0 ).
A mediana do Brasil como um todo foi 43,1, sendo a maior mediana obtida na região Nordeste ( 46,3 ), e a menor obtida na Norte ( 39,9 ).
A menor nota mínima foi 11,3, obtida na região Sudeste, e a maior nota mínima foi 20,6, obtida na região Nordeste.
O comportamento das notas dos estudantes de todo o Brasil pode ser observado no Gráfico 3.1 que apresenta um histograma com a distribuição das mesmas.
Essa é uma distribuição unimodal com moda no intervalo ( 40 ; 50 ].
Apesar do coeficiente de assimetria da distribuição das notas ser positivo ( 0,004 ), é tão pequeno que podemos considerar que a distribuição é simétrica.
As distribuições por Grande Região também apresentam coeficiente de assimetria positivo, com primeiro algarismo significativo na \textbf{ordem} dos centésimos, portanto, são aproximadamente simétricas.
A única exceção é a região Sudeste, com coeficiente de assimetria negativo ( – 0,11 ), indicando pequena concentração à direita e cauda um pouco mais longa à esquerda.
Os Gráficos 3.2, 3.3 e 3.4 apresentam informações referentes à média da nota final dos Participantes, desagregando os resultados de acordo com, respectivamente, as Grandes Regiões do país, a Categoria Administrativa e a Organização Acadêmica.
Os gráficos apresentam o valor da média das notas como uma barra e os extremos do intervalo de confiança de 95 \% como linhas verticais unidas por uma linha horizontal na forma da letra H maiúscula.
Considerando -se o gráfico de notas segundo Grande Região ( Gráfico 3.2 ), observa -se que não existe diferença estatisticamente significativa ao nível de 95 \% entre as médias.
Levando -se em conta os estudantes por Categorias Administrativas ( Gráfico 3.3 ), observa -se que também não existe diferença estatisticamente significativa entre as médias das notas das IES Públicas ( 41,6 ) e Privadas ( 43,6 ).
Tendo como foco o Gráfico 3.4, que apresenta as notas médias das provas segundo Organização Acadêmica, contata -se que existe diferença estatisticamente significativa ao nível de 95 \% entre a maior média, dos estudantes provenientes de Centro Universitários ( 55,3 ), e as obtidas pelos alunos de Universidades ( 41,8 ) e Faculdades ( 43,6 ). 3.1.2 Estatísticas Básicas no Componente de Formação Geral A Tabela 3.2 apresenta as estatísticas básicas em relação ao componente da prova que avalia a Formação Geral dos estudantes concluintes.
Os alunos de todo Brasil obtiveram desempenho médio de 50,6.
A maior média foi obtida na região Nordeste ( 58,6 ), e a menor, nas regiões Sul ( 46,0 ).
As demais médias foram : 48,6 na região Norte ; 51,7 na região Nordeste ; e 51,0 na região Centro-Oeste.
Quanto à variabilidade, o desvio padrão das notas dos estudantes do Brasil como um todo foi 16,4.
Já o maior desvio padrão foi obtido na região Norte ( 19,8 ), e o menor, na região Nordeste ( 13,1 ).
Os demais desvios padrões foram : 15,4 na região Sudeste ; 17,3 na região Sul e 14,9 na região Centro-Oeste.
A maior nota no componente de Formação Geral da prova do ENADE foi obtida por pelo menos um aluno da região Sudeste ( 91,5 ), enquanto que a menor nota máxima foi obtida na região Norte ( 77,0 ).
Nas outras regiões as notas máximas foram : 80,0 na região Nordeste ; 83,5 na região Sul ; e 79,0 na região Centro-Oeste.
A mediana do Brasil como um todo foi 52,5, sendo a menor mediana encontrada na região Sul ( 48,3 ) e a maior encontrada O Gráfico 3.5 propicia a avaliação do desempenho dos estudantes no componente de Formação Geral a partir do histograma da distribuição das notas correspondentes.
A distribuição é unimodal, com moda em ( 50 ; 60 ], enquanto na prova como um todo a moda foi alcançada no intervalo ( 40 ; 50 ].
Nota -se, ainda, que no gráfico 3.5 as notas apresentam uma maior dispersão do que no Gráfico 3.1 ( distribuição das notas da prova ), o que pode ser confirmado pela comparação dos desvios padrões : 12,2 para a nota da prova como um todo e 16,4 para o componente de Formação Geral.
Para o componente de Formação Geral, o coeficiente de assimetria da distribuição das notas dos estudantes é negativo ( – 0,37 ), diferentemente da prova como um todo.
Observa -se no histograma uma leve concentração à direita e cauda maior à esquerda.
Em todas as Grandes Regiões os histogramas também possuem assimetria negativa, variando entre – 0,53 ( Centro-Oeste ) e -0,20 ( Nordeste ). na região Nordeste ( 58,5 ).
A nota mínima nesta parte foi zero nas regiões Norte e Sul.
Nas demais regiões as notas mínimas foram : 33,0 na região Nordeste ; 7,5 na região Sudeste ; e 15,0 na região Centro-Oeste.
O Gráfico 3.5 propicia a avaliação do desempenho dos estudantes no componente de Formação Geral a partir do histograma da distribuição das notas correspondentes.
A distribuição é unimodal, com moda em ( 50 ; 60 ], enquanto na prova como um todo a moda foi alcançada no intervalo ( 40 ; 50 ].
Nota -se, ainda, que no gráfico 3.5 as notas apresentam uma maior dispersão do que no Gráfico 3.1 ( distribuição das notas da prova ), o que pode ser confirmado pela comparação dos desvios padrões : 12,2 para a nota da prova como um todo e 16,4 para o componente de Formação Geral.
Para o componente de Formação Geral, o coeficiente de assimetria da distribuição das notas dos estudantes é negativo ( – 0,37 ), diferentemente da prova como um todo.
Observa -se no histograma uma leve concentração à direita e cauda maior à esquerda.
Em todas as Grandes Regiões os histogramas também possuem assimetria negativa, variando entre – 0,53 ( Centro-Oeste ) e -0,20 ( Nordeste ).
Nos Gráficos 3.6, 3.7 e 3.8 são apresentadas as informações referentes ao desempenho dos Concluintes no componente de Formação Geral, em diferentes agregações : Grande Região do país, Categoria Administrativa e Organização Acadêmica.
Observa -se pelo Gráfico 3.6 que existe diferença estatisticamente significativa entre a maior média das notas no Componente de Formação Geral, obtida na região Nordeste ( 58,6 ) e as médias das regiões Sul ( 46,0 ) e Sudeste ( 51,7 ).
Vemos que o intervalo de confiança mais largo é o da região Norte ; já o intervalo mais estreito é observado na região Sudeste.
Este fato \textbf{está} relacionado, também, com o tamanho da população envolvida, menor na região Norte ( 31 presentes ) do que na Sudeste ( 164 presentes ).
No Gráfico 3.7, que representa as notas médias no Componente de Formação Geral segundo Categoria Administrativa do país, observa -se que não existe diferença estatisticamente significativa entre as médias dos concluintes das IES Públicas ( 49,2 ) e Privadas ( 51,6 ).
Considerando -se o tipo de Organização Acadêmica, nota -se, no Gráfico 3.8, uma diferença estatisticamente significativa entre a maior média e as demais.
Nos Centros Universitários a média ( 62,0 ) é mais elevada do que nas Universidades ( 49,8 ) e Faculdades ( 51,4 ). 3.1.3 Estatísticas Básicas do Componente de Conhecimento Específico A Tabela 3.3 apresenta as estatísticas básicas referentes ao componente de Conhecimento Específico da área de Tecnologia em Manutenção Industrial.
A média do desempenho dos alunos do Brasil como um todo foi 40,2.
A maior média foi obtida na região Sul ( 42,7 ), e a menor, na região Norte ( 36,6 ).
As demais médias foram : 42,1 na região Nordeste ; 38,7 na região Sudeste ; e 39,1 na região Centro-Oeste.
Quanto à variabilidade das notas, o desvio padrão do Brasil como um todo foi 13,4, sendo o maior desvio padrão observado na região Nordeste ( 14,2 ) e o menor na região Norte ( 11,1 ).
Os demais desvios foram : 13,4 da região Sudeste ; 13,1 da região Sul ; e 14,0 na região Centro-Oeste.
A mediana das notas dos estudantes de todo o Brasil foi 39,3.
A menor mediana ocorreu na região Centro-Oeste ( 37,5 ) e a maior na região Sul ( 42,7 ).
As demais medianas foram : 38,2 na região Norte, 41,7 na região Nordeste e 38,0 na região Sudeste.
A nota máxima do Brasil como um todo foi 85,0, sendo obtida por pelo menos um aluno da região Sul.
As demais notas máximas foram : 61,0 na região Norte ; 76,0 na região Nordeste ; 69,8 na região Sudeste ; e 80,5 na região Centro-Oeste.
A nota mínima zero ocorreu apenas na região Sudeste.
A maior nota mínima ocorreu na região Sul ( 12,8 ).
Assim como os Gráficos 3.1 e 3.5, o Gráfico 3.9, apresentado a seguir, proporciona uma avaliação do desempenho de concluintes em relação ao componente de Conhecimento Específico com um histograma da distribuição das notas correspondentes.
Esta também é uma distribuição unimodal, e o intervalo modal é o ( 30 ; 40 ].
O coeficiente de assimetria da distribuição das notas do componente de Conhecimento Específico é positivo e bem próximo de zero ( 0,10 ).
Nota -se pelo histograma ( Gráfico 3.9 ) que esta é uma distribuição aproximadamente simétrica.
Já na região Norte, o coeficiente de assimetria é negativo, igual a – 0,50, evidenciando que a distribuição das notas dos estudantes dessa região, no componente de Conhecimento Específico, tem cauda maior à esquerda.
Os Gráficos 3.10, 3.11 e 3.12 apresentam uma comparação dos resultados em relação à Grande Região do país, à Categoria Administrativa e à Organização Acadêmica, agora levando em conta o desempenho de alunos no componente de Conhecimento Específico da prova.
Pelo Gráfico 3.10, observa -se que não existe diferença estatisticamente significativa entre as médias das notas no componente de Conhecimento Específico das Grandes Regiões.
Quanto à Categoria Administrativa ( gráfico 3.11 ), observa -se um comportamento semelhante àquele da parte de Formação Geral e à prova como um todo, ou seja, não existe diferença estatisticamente significativa entre as médias das IES Públicas e Privadas, sendo que a menor média foi obtida por alunos de IES Públicas de ensino.
Quanto ao Gráfico 3.12, mais uma vez, observa -se que existe diferença estatisticamente significativa ao nível de 95 \% das notas dos Concluintes de Centro Universitários no componente de Conhecimento Específico em relação aos demais tipos de Organização Acadêmica.
Sendo que a média dos Concluintes de Centros Universitários ( 53,1 ) foi maior do que de Universidades ( 39,2 ) e de Faculdades ( 41,0 ). 3.2 ANÁLISE DAS QUESTÕES OBJETIVAS 3.2.1 Componente de Formação Geral A Tabela 3.4 apresenta as estatísticas básicas relativas às oito questões objetivas do componente da prova que abrange a Formação Geral dos estudantes.
A média do Brasil foi 50,0.
A maior média foi encontrada na região Nordeste ( 54,8 ) e a menor na região Norte ( 46,0 ).
As demais médias foram : 52,7 na região Sudeste ; 46,2 na região Sul ; e 47,6 na região Centro-Oeste.
O desvio padrão do Brasil foi 18,9, sendo o maior desvio padrão encontrado na região Norte ( 22,2 ), e o menor, na região Nordeste ( 17,2 ).
Os demais desvios foram : 17,7 na região Sudeste ; 19,5 na região Sul ; e 19,1 na região Centro-Oeste.
As medianas ( 50,0 ) foram iguais para todas as regiões.
A nota máxima 100,0 foi obtida por pelo menos um aluno da região Sudeste, nas demais regiões a nota máxima foi 87,5.
A nota mínima foi zero em todas as regiões, exceto na região Nordeste ( 12,5 ).
A Tabela 3.5 apresenta os índices de facilidade e discriminação ( ponto bisserial ) das questões objetivas do componente de Formação Geral.
Quanto ao índice de facilidade, foram usadas as seguintes cores para diferenciar o nível de dificuldade da questão : • Azul para as questões classificadas com índice muito fácil ( índice \&gt ; =0,86 ), verde para as questões classificadas com índice fácil ( 0,61 a 0,85 ), amarelo para as questões classificadas com médio ( 0,41 a 0,60 ), vermelho para as questões classificadas com difícil ( 0,16 a 0,40 ) e roxo para as questões classificadas com muito difícil ( \&lt ; =0,15 ).
Já quanto ao índice de discriminação, foram usadas as seguintes cores para qualificar a questão : • As questões classificadas com índice fraco receberam a cor vermelho ( índice \&lt ; =0,19 ), as classificadas com médio receberam a cor amarelo ( 0,20 a 0,29 ), as classificadas com bom receberam a cor verde ( 0,30 a 0,39 ) e as classificadas com muito bom ( \&gt ; =0,40 ) receberam a cor azul.
As questões objetivas do componente de Formação Geral, segundo o índice de facilidade, foram assim avaliadas : das oito questões, nenhuma teve o índice de facilidade classificado como muito fácil.
Quatro questões foram tidas como fáceis, por terem índice de acertos situado na faixa entre 0,61 e 0,85 ( de 61,0 \% a 85,0 \% de acertos ).
Uma questão foi considerada de dificuldade média, situando -se no intervalo entre 0,41 e 0,60 do índice de facilidade, ou seja, houve entre 41,0 \% e 60,0 \% de acertos.
Outras três questões foram classificadas na categoria difícil, situando -se no intervalo entre 0,16 e 0,40.
Por fim, nenhuma das questões apresentou menos de 15 \% de acertos, razão pela qual seriam classificadas como muito difíceis.
A Tabela 3.5 apresenta os índices de facilidade e discriminação ( ponto bisserial ) das questões objetivas do componente de Formação Geral.
Quanto ao índice de facilidade, foram usadas as seguintes cores para diferenciar o nível de dificuldade da questão : • Azul para as questões classificadas com índice muito fácil ( índice \&gt ; =0,86 ), verde para as questões classificadas com índice fácil ( 0,61 a 0,85 ), amarelo para as questões classificadas com médio ( 0,41 a 0,60 ), vermelho para as questões classificadas com difícil ( 0,16 a 0,40 ) e roxo para as questões classificadas com muito difícil ( \&lt ; =0,15 ).
Já quanto ao índice de discriminação, foram usadas as seguintes cores para qualificar a questão : • As questões classificadas com índice fraco receberam a cor vermelho ( índice \&lt ; =0,19 ), as classificadas com médio receberam a cor amarelo ( 0,20 a 0,29 ), as classificadas com bom receberam a cor verde ( 0,30 a 0,39 ) e as classificadas com muito bom ( \&gt ; =0,40 ) receberam a cor azul.
As questões objetivas do componente de Formação Geral, segundo o índice de facilidade, foram assim avaliadas : das oito questões, nenhuma teve o índice de facilidade classificado como muito fácil.
Quatro questões foram tidas como fáceis, por terem índice de acertos situado na faixa entre 0,61 e 0,85 ( de 61,0 \% a 85,0 \% de acertos ).
Uma questão foi considerada de dificuldade média, situando -se no intervalo entre 0,41 e 0,60 do índice de facilidade, ou seja, houve entre 41,0 \% e 60,0 \% de acertos.
Outras três questões foram classificadas na categoria difícil, situando -se no intervalo entre 0,16 e 0,40.
Por fim, nenhuma das questões apresentou menos de 15 \% de acertos, razão pela qual seriam classificadas como muito difíceis.
Como já comentado, para análise das questões objetivas relativas à Formação Geral segundo o poder de discriminação, utilizou -se, o índice de discriminação ponto bisserial.
Nesta análise as questões foram assim avaliadas : quatro das oito questões apresentaram índices acima ou igual a 0,40 e, assim, foram classificadas com índice muito bom para esse grupo de estudantes.
Três questões tiveram bom índice de discriminação, com valor entre 0,30 e 0,39 e uma outra apresentou um valor entre 0,20 e 0,29, sendo classificada com índice médio para esse grupo de estudantes.
Nenhuma questão teve nível fraco de discriminação.
O índice de facilidade variou de 0,19 a 0,75, e o de discriminação, de 0,27 a 0,53.
As quatro questões com índices de discriminação muito bom, as de números 1, 2, 3 e 5, figuraram entre os diversos níveis de dificuldade desse conjunto : três classificadas na categoria fácil ( questões 1, 3 e 5 ) do índice de facilidade e uma na categoria difícil ( questão 2 ).
Em particular, a questão 3 foi a que apresentou o maior poder discriminatório, com índice 0,53, e foi também uma das mais fáceis, com uma proporção de 0,64 acertos.
O máximo de acertos foi alcançado pela questão 6 com um índice de facilidade de 0,75.
A questão de número 4 apresentou um índice de facilidade de 0,51, ou seja, um pouco mais da metade dos estudantes conseguiu resolvê -la, dentro do universo de participantes.
Seu índice de discriminação foi bom ( 0,37 ).
Já a questão 8 obteve índice de discriminação médio, 0,27, e seu índice de facilidade foi difícil ( 0,19 ).
O Gráfico 3.13, para exemplificar, analisa o comportamento da questão de número 3 de Formação Geral.
Trata -se da segunda questão mais fácil e a que obteve o maior índice de discriminação dessa parte da prova.
Neste gráfico, \textbf{cada} uma das cinco curvas representa o percentual de respostas em determinada alternativa da questão, em função do número de acertos dos estudantes nesta parte da prova ( Formação Geral/Múltipla Escolha ), antes de possíveis eliminações pelo critério do ponto bisserial.
Em princípio, a soma das escolhas possíveis deveria ser igual a 100 \%.
Não é este o caso, notadamente para o número zero de acertos, pois não aparecem no gráfico as questões deixadas em branco ou com múltiplas respostas.
Como foram oito as questões, os valores variam de 0 a 8 acertos.
A curva em vermelho corresponde à alternativa E, a correta para esta questão.
Assim, observa -se que entre os estudantes com menor número de acertos, nessa parte do exame, a situação mais frequente foi a escolha de uma das alternativas incorretas : a alternativa A ( em azul ) ou C ( em preto ).
Na medida em que o número de acertos aumenta, indicando desempenho melhor nesta parte da prova, aumenta concomitantemente a proporção de estudantes que selecionaram a alternativa correta E, atingindo 100 \% para os estudantes com 8 acertos.
Essa análise permite verificar como a questão discriminou os grupos de desempenho, justificando o alto índice obtido na questão.
Os gráficos relativos às demais questões de Formação Geral constam do Anexo I. 3.2.2 Componente de Conhecimento Específico A Tabela 3.6 apresenta as estatísticas básicas em relação às questões objetivas do componente de Conhecimento Específico da prova, por Grande Região.
A média do Brasil deste componente foi de 43,2.
A menor média foi observada na região Norte ( 39,5 ) e a maior na região Sul ( 46,4 ).
O desvio padrão de todo o Brasil foi 14,6, sendo o menor desvio padrão encontrado na região Norte ( 12,1 ), e o maior, na região Nordeste ( 15,6 ).
A mediana de todo o Brasil foi 45,0, a mesma encontrada nas regiões Nordeste e Sul.
Nas demais regiões do Brasil a mediana foi 40,0.
A nota máxima da prova ( 85,0 ) foi obtida no componente de Conhecimento Específico por pelo menos um aluno da região Sul.
A menor nota máxima ocorreu na região Norte ( 65,0 ).
As notas mínimas foram : 0,0 na região Sudeste ; 5,0 nas regiões Norte e Nordeste ; 10,0 na região Centro-Oeste ; e 15,0 na região Sul.
A Tabela 3.7 apresenta os índices de facilidade e discriminação ( ponto bisserial ) das questões objetivas do componente de Conhecimento Específico.
Para facilitar a diferenciação das questões usaram -se as mesmas cores da Tabela 3.5 para as diferentes classificações dos índices de facilidade e de discriminação.
Dentre as questões objetivas da parte da prova relativa ao componente de Conhecimento Específico, a de número 28 foi anulada pela Comissão.
Desse modo, a classificação quanto ao índice de facilidade foi estabelecida com base em 26 das 27 questões.
A partir dos índices obtidos, pode -se concluir que a maioria das questões objetivas da prova foi considerada pelo difícil ( a categoria modal ) : das 26 questões, quinze foram classificadas como difíceis.
Não houve questão classificada como muito fácil, nem tampouco classificada de muito difícil, ao passo que três foram tidas como fáceis, na faixa de 0,61 a 0,85 do índice de facilidade, e oito consideradas médias, entre 0,41 e 0,60.
Já quanto aos índices de discriminação das questões objetivas do componente de Conhecimento Específico da prova, tem -se como resultado a seguinte classificação : oito das 26 questões válidas foram consideradas como boas, enquanto somente uma delas teve índice de discriminação muito bom.
Assim, para cerca de um terço das questões – nove em 26 – os índices de discriminação foram bons ou muito bons.
Dentre as demais, onze delas foram classificadas como médias e outras seis como fracas, sendo dezessete, por conseguinte, a quantidade de questões nos dois patamares mais baixos de discriminação.
Constata -se, assim, que a prova – no que se refere ao componente de Conhecimento Específico – possuía capacidade baixa de discriminar entre aqueles que dominam ou não o conteúdo.
A questão que alcançou o maior índice de discriminação, a de número 23, classificada com índice muito bom, situando -se no intervalo de 0,40 a 0,45 do índice, foi classificada na categoria médio quanto ao índice de facilidade.
As questões de número 11 e 13 foram as mais difíceis dentre as 26 questões específicas válidas, com baixo índice de facilidade, apenas 17,0 \% de acertos.
A primeira destas questões, a de número 11, apresentou poder discriminatório igualmente baixo, 0,07.
Já a questão 13 apresentou um índice de discriminação médio, 0,23.
Destaca -se, também, a questão 10, com índice de facilidade 0,21, o que, em termos percentuais, corresponde a 21,0 \% de estudantes que responderam acertadamente, obtendo, ainda, 0,12 de índice de discriminação.
No total, foram seis as questões eliminadas pelo critério do ponto bisserial : 10, 11, 18, 19, 34 e 35 e não foram utilizadas no cômputo da nota final.
A título de exemplo das análises do comportamento das questões objetivas, o Gráfico 3.14 analisa a questão 23 do componente de Conhecimento Específico.
Esta foi uma questão de nível de facilidade médio, apresentando índice de facilidade 0,42, ou seja, 42,0 \% dos estudantes assinalaram acertadamente a opção E, correspondente ao gabarito.
Seu índice de discriminação foi igual a 0,44, classificado como muito bom.
Neste gráfico, \textbf{cada} uma das cinco curvas representa o percentual de respostas em determinada alternativa da questão 23, em função da nota dos estudantes nesta parte da prova, antes de possíveis eliminações de questões pelo critério do ponto bisserial.
A alternativa correta E, representada no gráfico pela curva em vermelho, foi escolhida em maiores proporções pelos alunos com desempenho melhor nesta parte da prova.
Já as alternativas incorretas, também denominadas distratores, foram selecionadas principalmente por aqueles com notas mais baixas.
Aqueles com nota zero, na sua quase totalidade deixaram esta questão em branco ou marcaram mais de uma alternativa, comportamento considerado inválido.
A proporção de alunos que selecionou a resposta correta E aumenta gradativamente, chegando a atingir 100 \% para 17 acertos ou mais, enquanto a proporção dos que escolheram alternativas incorretas decai, a partir da primeira nota não nula, como função da nota nesta parte da prova.
Os gráficos relativos às demais questões do componente de Conhecimento Específico constam do Anexo I. 3.3 ANÁLISE DAS QUESTÕES DISCURSIVAS 3.3.1 Componente de Formação Geral As análises dos resultados de desempenho dos estudantes de Tecnologia em Manutenção Industrial nas duas questões discursivas relativas à Formação Geral encontram -se na Tabela 3.8 e no Gráfico 3.15.
Na tabela 3.8 observa -se que as notas médias foram um pouco mais altas nesse conjunto de questões do que no das objetivas.
Os estudantes de todo o Brasil obtiveram, em Formação Geral, média 50,0 nas questões objetivas e 51,6 nas questões discursivas.
A mediana confirma o desempenho um pouco melhor dos alunos de todo o Brasil nas questões discursivas do componente de Formação Geral.
Enquanto essa estatística foi de 53,8 para questões discursivas, para as questões objetivas essa estatística foi 50,0.
Pode -se, também, notar um aumento do desvio padrão de 18,9, nas questões objetivas do componente de Formação Geral dos alunos de todo o Brasil, para 24,9 nas questões discursivas do mesmo componente.
A maior mediana encontrada na região Nordeste ( 68,8 ) e a menor na região Sul ( 45,0 ).
A nota máxima ( 100,0 ) foi obtida por pelo menos um aluno da região Sudeste e a nota mínima foi zero para todas as regiões, sem exceção.
O Gráfico 3.15 representa a distribuição das notas nas questões discursivas no componente de Formação Geral.
A moda desta distribuição ocorre no intervalo, ( 60 ; 70 ].
O coeficiente de assimetria da distribuição de notas de todo o Brasil é negativo ( – 0,37 ) ( assimetria à esquerda ).
Apenas na região Sul o comportamento da distribuição é aproximadamente simétrico, com coeficiente positivo e próximo de zero ( 0,05 ).
Na sequência, os resultados verificados para \textbf{cada} uma das questões discursivas de Formação Geral serão apresentados, estabelecendo -se relações com os conteúdos abordados em \textbf{cada} uma delas.
Os comentários da Banca de docentes corretores a respeito do observado na correção das respostas dos estudantes, suas impressões e conclusões serão apresentados junto à análise de \textbf{cada} questão.
Cumpre esclarecer que, tendo em vista que as questões discursivas de Formação Geral são padronizadas, ou seja, constam de todas as provas, os comentários da Banca são os mesmos para todas as carreiras acadêmicas, sendo direcionados a todos os estudantes que participaram do ENADE/2011.
A seguir, serão analisados os desempenhos da Área de Tecnologia em Manutenção Industrial nas duas questões discursivas de Formação Geral do ENADE/2011, comparando os resultados obtidos com comentários para \textbf{cada} questão. 3.3.1.1 Análise da Questão Discursiva 1 do Componente de Formação Geral Os dados de Tecnologia em Manutenção Industrial, obtidos a partir das respostas à questão 1, encontram -se na Tabela 3.9 e no Gráfico 3.16.
Nessa questão – de melhor desempenho dentre as duas de Formação Geral – os alunos de todo Brasil tiveram média 58,5.
A maior média para a questão 1 foi obtida na região Nordeste ( 69,9 ), e a menor, na região Sul ( 52,4 ).
Quanto à variabilidade das notas, o desvio padrão de todo o Brasil foi 32,2.
O menor desvio padrão foi obtido na região Nordeste ( 28,5 ), e o maior desvio padrão ocorreu na região Centro-Oeste ( 33,2 ).
As notas máximas e mínimas da questão discursiva 1 foram as mesmas para todas as regiões do Brasil, respectivamente, 100,0 e 0,0.
A mediana de todo o Brasil foi 60,0, a mesma da região Sudeste.
A menor mediana ocorreu na região Sul ( 50,0 ), e a maior, nas regiões Nordeste e Centro-Oeste ( 70,0 ).
O Gráfico 3.16 mostra a distribuição das notas na questão discursiva 1 do componente de Formação Geral.
Observa -se que a moda é o intervalo ( 90 ; 100 ] e que outros máximos locais ocorrem nos intervalos ( 20 ; 30 ], ( 60 ; 70 ] e ( 80 ; 90 ].
O coeficiente de assimetria é negativo ( – 0,27 ), o que explica a cauda mais longa à esquerda.
Em todas as Grandes Regiões o comportamento das distribuições possui assimetria negativa.
Os coeficientes de assimetria variam entre – 0,67 ( Nordeste ) a – 0,02 ( Sul ). 3.3.1.2 Comentários sobre a correção das respostas à Questão Discursiva 1 De maneira geral, a aplicação da chave de correção da questão 1 de Formação Geral não apresentou qualquer dificuldade digna de menção.
Para isso contribuíram, indubitavelmente, os ajustes feitos após a correção da amostra e a reunião entre todos os membros da Banca de docentes corretores.
As poucas dúvidas, todas pontuais, apresentadas pelos corretores, foram acompanhadas e respondidas pela coordenação e subcoordenação da correção das questões de Formação Geral, por meio da ferramenta de Gerenciamento de Dúvidas do Sistema de Correção On-line.
Não houve registro de qualquer ocorrência que pusesse em xeque o padrão de resposta ou a efetividade e a adequação da chave de correção.
Explica -se : trata -se de questão com comando claro, direto e objetivo ( solicitava -se, basicamente, três vantagens justificadas de cursos a distância ), cujas respostas foram corrigidas por meio da aplicação de um chave de correção testada e aprovada previamente.
Havia absoluta clareza quanto aos critérios de avaliação da correspondência entre as respostas dos estudantes e as possibilidades de vantagens de cursos a distância admitidas como corretas no padrão de resposta oficial, além de gradações explícitas ( e fáceis de aplicar ) dos diferentes níveis de pontuação previstos.
Felizmente, portanto, não há reparo a registrar em relação à facilidade de aplicação do padrão de resposta e da chave de correção, e nem em relação à atribuição dos diferentes níveis de pontuação previstos.
Todas as dificuldades que poderiam ter obstado a correta aplicação do padrão de resposta oficial e da respectiva chave de correção foram evitadas por meio dos ajustes feitos após a correção da amostra e fartamente debatidos com toda a Banca.
Digno de nota é que quantidade tão significativa de profissionais envolvidos na mesma tarefa – tanto para a questão 1 quanto para a 2 – tenha apresentado tão poucas dificuldades na execução da correção, em termos proporcionais.
Em suma, a correção da questão 1 da prova de Formação Geral do ENADE 2011 foi exemplarmente bem planejada, servindo -se de padrão de resposta muito bem adequado à questão proposta.
Quanto ao tema desta questão, em particular, a Banca verificou que uma parcela significativa de estudantes evocou experiências bastante concretas e próximas de sua realidade.
Houve várias respostas que indicavam uma vivência pessoal de ensino superior na modalidade Educação a Distância ( EaD ), evidenciando o tom de depoentes nos textos apresentados.
Os estudantes, em sua grande maioria, utilizaram parte considerável do espaço de 15 linhas disponíveis para a resposta – e outra parcela menos significativa dos que não o fizeram demonstraram notável capacidade de atender ao comando da questão de maneira objetiva, curta e, via de regra, correta.
Registre -se, ainda, que foram relativamente poucos os casos de respostas que tenham passado ao largo do tema em pauta na questão.
As capacidades de leitura, de compreensão do comando proposto e de expressão escrita que os estudantes avaliados na edição 2011 do ENADE foram satisfatórias.
Não obstante, seria leviano perder de vista que a qualidade dos textos redigidos em resposta às questões discursivas do Exame ainda \textbf{está} muito aquém do que se espera de concluintes de cursos de ensino superior de todas as regiões do país.
Quanto ao conteúdo das respostas, a Banca constatou boa capacidade, por parte da maioria dos estudantes, de compreensão do tema e do comando da questão.
Foram relativamente poucos os casos de respostas que deixaram de enumerar vantagens da modalidade EaD, e proporcionalmente escassos os estudantes que citaram vantagens não previstas no padrão de resposta.
Foi frequente, a tentativa direta de atender ao comando da questão.
Os erros mais comuns, em relação ao padrão de respostas e à grade de correção, foram fruto do desdobramento em vários “ itens ” daquilo que, de acordo com o padrão de respostas oficial, representava uma única vantagem.
Destacaram -se, neste caso, as respostas que apontavam a flexibilidade de horário e/ou local como duas vantagens distintas daquela modalidade de ensino.
Quanto aos diferentes níveis de pontuação previstos, a maior causa de baixas pontuações foi a ausência de justificativas, e mesmo de argumentação, para uma ou mais das vantagens enumeradas.
Isso demonstra não apenas a objetividade das respostas, coerente com a objetividade do comando da questão ( “ enumere três vantagens de um curso a distância ” ), mas também \textbf{certa} dificuldade de formulação plena de um texto, ou ao menos de parágrafos, em formato dissertativo – mesmo diante de uma média de 5 linhas disponíveis para \textbf{cada} vantagem a enumerar.
Os acertos mais comuns, ou seja, os “ itens ” do padrão de respostas mais frequentemente mencionados foram : ( 1 ) a flexibilidade de horário e/ou local ; ( 2 ) a capilaridade do ensino a distância ; ( 3 ) a democratização do acesso à educação de qualidade ; e ( 4 ) os custos menores que os de cursos presenciais.
Dentre as vantagens previstas no padrão de respostas que foram menos citadas, destacam -se a inclusão de pessoas com comprometimento motor, a qualificação de professores e a troca de experiências entre os participantes.
Foram muito comuns, no entanto, as menções à supostamente maior facilidade de acesso a professores ou/e tutores em cursos superiores a distância. 3.3.1.3 Análise da Questão Discursiva 2 do Componente de Formação Geral A Tabela 3.10 mostra que o desempenho dos estudantes na questão 2 ( média 44,7 ), inferior ao obtido na questão de número 1 ( média 58,5 ).
A região Nordeste foi aquela onde a média, nessa questão, foi maior ( 58,6 ), e a de menor média foi a região Sul ( 38,9 ).
Quanto à variabilidade das notas, o desvio padrão de todo o Brasil foi 28,9, inferior ao obtido na questão de número 1 ( 32,2 ).
O maior desvio nessa questão foi obtido na região Norte ( 29,9 ), enquanto o menor foi obtido na região Nordeste ( 24,0 ).
A mediana comprova a diferença de desempenho entre essas duas questões : enquanto a mediana da questão 1 foi 60,0, na questão 2 foi 50,0.
Nesta questão, a maior mediana ocorreu na região Nordeste ( 65,0 ) e a menor na região Sul ( 40,0 ).
A nota máxima ( 100,0 ) foi obtida em todas as regiões do Brasil, exceto na região Norte, onde a nota máxima foi 80,0.
A nota mínima na questão 2 foi zero em todas as regiões.
O Gráfico 3.17 mostra a distribuição das notas na questão discursiva 2 do componente de Formação Geral.
O intervalo modal, o de maior frequência, é aquele que representa as questões deixadas em branco.
Desconsiderando -se este intervalo, observa -se um ponto de máximo no intervalo ( 50 ; 60 ].
Com exceção da região Sul ( 0,14 ), os coeficientes de assimetria de todas as regiões são negativos, sendo a região Nordeste a que possui assimetria esquerda mais forte ( – 1,06 ).
Para a distribuição de todas as notas do Brasil o coeficiente de assimetria é igual a – 0,30. 3.3.1.4 Comentários sobre a correção das respostas à Questão Discursiva 2 Diferentemente da questão 1, a questão 2, cujo tema aborda políticas públicas para a erradicação do analfabetismo e para o aumento da empregabilidade, permitiu mais amplitude na elaboração das respostas dos estudantes, uma vez que solicita a apresentação de uma proposta.
À falta de objetividade técnica na elaboração da questão – sobretudo no seu comando – correspondeu uma miríade de possibilidades interpretativas.
Registre -se que não raras foram as respostas que se utilizaram de lugares comuns e exortações religiosas/humanitárias/cívicas, contudo, na grande maioria dos casos, essas exortações foram usadas como um encerramento do texto e não comprometeram a resposta.
Muito frequente também foi a confusão entre política pública e política partidária, bem como alguns poucos, confundiam programas educacionais com programas educativos nas TVs.
Assim, destacaram -se como propostas/programas mais recorrentes :  prosseguimento das políticas já existentes, mas com o aumento dos investimentos, normalmente sugerindo a ampliação da bolsa família, da Educação de Jovens e Adultos ( EJA ) e da EaD ;  valorização do magistério, construção de escolas e melhoria das atuais ;  parceria do governo com empresas para manutenção de salas de aula visando a alfabetização de seus funcionários, oferecimento de estágios e redução da carga horária em troca de isenção de impostos ;  parcerias com igrejas e ONGs para criação de espaços de alfabetização ;  escolas itinerantes e alfabetizadores em domicílio, principalmente para pessoas com dificuldade de locomoção, como os idosos e deficientes físicos ;  erradicação do trabalho infantil ;  vinculação da bolsa família não apenas à frequência, mas também e, principalmente, ao resultado obtido pelo aluno na escola ;  revisão das políticas atuais, sendo a mais frequente a extinção da aprovação automática ;  críticas consistentes em relação a modelos didáticos considerados inadequados e desestimulantes para a educação não só de adultos, mas de pessoas de todas as idades.
Destacam -se, ainda, outras sugestões apresentadas :  Creches nas escolas onde os pais \textbf{estão} sendo alfabetizados.  Diminuição de duas horas na jornada de trabalho em empresas para funcionários não alfabetizados, para que possam frequentar a escola.  Campanhas educativas vinculadas aos meios de comunicação.  Aumento do número de escolas noturnas.  Formação específica para professores alfabetizadores.
Quanto à relação entre o analfabetismo e a empregabilidade, deve -se sublinhar que nem todos os estudantes estabeleceram claramente o vínculo entre essas duas situações sociais.
Alguns falaram separadamente de uma e de outra.
Mas a maioria fez referência à necessidade de estudo para “ conseguir um bom emprego com um bom salário ”.
Alguns, em menor número, estabeleceram de forma bastante interessante a questão histórica para a situação do Nordeste ; e, também, a relação entre a pessoa analfabeta/com pouca instrução e a desempregada/empregada em serviços mais pesados e pior remunerados, que não teriam condições de educar seus filhos que, por sua vez, também não teriam melhores oportunidades no mercado de trabalho, identificando a formação de um círculo vicioso e a necessidade de sua interrupção.
Muitos afirmaram que o analfabetismo não é o único responsável pelo desemprego e sim a má distribuição de renda.
Grande também foi o número de alunos que criticaram o resultado da pesquisa, afirmando que a mesma é enganosa, uma vez que considera alfabetizada a pessoa que “ desenha seu nome ”, sendo comum associarem a atual política de alfabetização com ganhos eleitorais.
Foi comum a resposta incompleta, em que o estudante fez apenas a análise das desigualdades/crítica do quadro apresentado, ou só apresentou proposta.
Alguns se limitaram a responsabilizar o governo referindo -se de forma bastante genérica à questão da “ educação ” e “ profissional ”.
Entre os equívocos que mais se repetiram, destacou -se a simples análise dos dados apresentados na tabela que consta do enunciado.
Alguns textos, inclusive, \textbf{estavam} corretos, sem que, no entanto, fosse respondida a pergunta.
Da mesma forma, quando os estudantes partiram para a segunda parte da questão, fizeram referência a vários programas já existentes ou simplesmente disseram que é muito importante que existam projetos para a educação.
Outro equívoco recorrente foi a análise da educação no Brasil como um todo.
Essa análise, apesar de correta, não se referia especificamente à questão do analfabetismo, o que, por conseguinte, levava a sugestões que não eram direcionadas à erradicação do analfabetismo, e sim à melhora da educação no Brasil.
Sendo assim, foram apresentadas propostas como o aumento do número de faculdades, o reforço de alunos do Ensino Médio, a criação de escolas técnicas, etc.
Por outro lado, houve um grande número de redações bem escritas e precisas, no que se refere ao que foi exigido pela questão.
Foram análises equilibradas e sensatas, correlacionando a problemática do analfabetismo com o desemprego e a apresentação de sugestões bastante consistentes. 3.3.2 Componente de Conhecimento Específico Na parte da prova relativa às questões discursivas no componente de Conhecimento Específico ( Tabela 3.11 ), observa -se que a média foi mais baixa do que para as questões discursivas do componente de Formação Geral ( Tabela 3.8 ).
Enquanto no componente de Formação Geral a média para estudantes de Tecnologia em Manutenção Industrial de todo o Brasil foi 51,6, na parte de Conhecimento Específico a média foi 23,1.
A maior média deste componente foi obtida pelos estudantes da região Nordeste ( 28,0 ), e a menor, pelos da região Norte ( 19,7 ).
Quanto à variabilidade das notas, o desvio padrão de todo o Brasil foi 15,4.
O maior desvio padrão foi encontrado na região Sudeste ( 16,0 ), e o menor, na região Norte ( 13,2 ).
A maior nota máxima foi obtida na região Sul ( 85,0 ), enquanto a menor nota máxima foi encontrada na região Norte ( 40,0 ).
Além disso, a nota mínima ( 0,0 ) foi obtida por alunos de todas as regiões do Brasil, sem exceção.
A mediana do Brasil como um todo foi 23,3.
A maior mediana foi a da região Nordeste ( 28,3 ), e a menor ( 21,7 ) ocorreu nas regiões Norte, Sudeste e Centro-Oeste.
O Gráfico 3.18 representa a distribuição das notas nas questões discursivas no componente de Conhecimento Específico.
Essa distribuição é unimodal com moda no intervalo de ( 20 ; 30 ].
Nota -se uma tendência decrescente a partir do intervalo modal.
Existe um segundo máximo local no primeiro intervalo, [ 0 ; 10 ].
Não foram registradas notas nos intervalos ( 70 ; 80 ] e ( 90 ; 100 ].
O coeficiente de assimetria é positivo e pequeno ( 0,48 ).
Este coeficiente também é positivo nas regiões Sudeste, Sul e Centro-Oeste.
O histograma apresenta leve assimetria à esquerda ( coeficiente negativo ) na regiões Norte e Nordeste. 3.3.2.1 Análise da Questão Discursiva 3 do Componente de Conhecimento Específico Na questão 3, cujos resultados aferidos encontram -se descritos na Tabela 3.12, a média dos estudantes de todo o Brasil foi 20,9.
A menor média nessa questão foi obtida pelos alunos da região Centro-Oeste ( 15,5 ), enquanto a maior média foi obtida na região Sudeste ( 22,0 ).
Quanto à variabilidade das notas, o desvio padrão de todo o Brasil foi 22,3.
O maior desvio padrão foi obtido na região Sudeste ( 24,4 ), enquanto o menor foi obtido na região Centro-Oeste ( 17,0 ).
A nota máxima, 100,0 pontos, foi alcançada por pelo menos um aluno da região Sudeste.
A região onde ocorreu a menor nota máxima foi a Norte ( 70,0 ).
A mediana do Brasil como um todo foi 15,0, a mesma nas regiões Sudeste, Sul e Centro-Oeste.
Nas demais regiões a mediana foi 20,0.
A nota mínima ( 0,0 ) foi obtida em todas as regiões do Brasil, sem exceção.
O Gráfico 3.19 mostra a distribuição das notas na questão discursiva 3, do componente de Conhecimento Específico.
Essa distribuição a moda ocorre na categoria “ em branco ”.
Dentre os alunos que tentaram resolver a questão, cerca de 70 \% deles, a maior freqüência de notas ocorreu no intervalo ( 20 ; 30 ].
A distribuição tem forte assimetria positiva, concentração do lado esquerdo e cauda longa do lado direito ( coeficiente de assimetria 1,20 ).
Em todas as regiões o comportamento da distribuição é assimétrico positivamente, com coeficientes variando de 0,74 ( Norte ) a 1,57 ( Centro-Oeste ). 3.3.2.2 Comentários sobre a correção das respostas à Questão Discursiva 3 Esta questão foi considerada, pela Banca de Docentes Corretores, como fácil.
Busca o conhecimento do discente quanto a um assunto muito comum no dia a dia da manutenção industrial. \textbf{Está}, portanto, dentro do contexto de avaliação para o nível de Tecnólogo de Manutenção Industrial.
Deixa entender ao discente que ele é a pessoa qualificada a executar o serviço.
Quanto ao desempenho dos estudantes, aproximadamente metade dos que responderam a questão não mencionaram o uso de EPIs.
Uma quantidade relativamente alta de concluintes não mencionaram o tipo de aparelho de medição a ser usado ou mencionaram aparelho inadequado para este tipo de medição, como multímetro ou amperímetro.
A grande maioria dos discentes não detalhou o procedimento usual de medição de resistência ao isolamento.
Alguns apenas mencionaram de forma geral.
No que se refere aos critérios e procedimentos de correção adotados para avaliar, foi adotado exatamente o solicitado pela Comissão Assessora da Área.
Buscou -se considerar as respostas dos estudantes em sua integridade, pontuando tudo o que \textbf{estivesse} conceitualmente correto, mesmo que o solicitado na questão não \textbf{estivesse} colocado na forma gramatical padrão.
Foi levada em consideração a ideia que o discente pretendeu colocar no texto.
Embora a questão 3 possa ser considerada fácil, o desempenho dos discentes foi de mediano para fraco, haja vista o número expressivo de participantes que não responderam a questão, deixando -a em branco.
Além disso, pelo pequeno número de alunos, em relação ao total, que obteve nota igual ou acima de 5,0.
Alguns alunos escreveram que não tiveram aula relacionada ao conteúdo da questão 3 no curso.
Além disso, dois estudantes mencionaram que sua especialidade é Manutenção Mecânica, e que, por isso, não conheciam o procedimento de medição de isolamento de motor. 3.3.2.3 Análise da Questão Discursiva 4 do Componente de Conhecimento Específico A Tabela 3.13 contém as informações relativas à questão 4 do conjunto de questões do componente de Conhecimento Específico.
A média do desempenho dos estudantes de todo o Brasil nesta questão foi 27,2, sendo a menor média registrada na região Norte ( 20,2 ) e a maior na região Nordeste ( 33,8 ).
A nota máxima ( 100,0 ) foi atingida por pelo menos um aluno da região Centro-Oeste.
A mediana das notas de todo o Brasil foi 30,0, a mesma das regiões Sudeste e Sul.
A nota mínima ( 0,0 ) foi obtida em todas as regiões do Brasil, sem exceção.
O Gráfico 3.20, representa a distribuição de notas na questão discursiva 4, no componente de Conhecimento Específico.
O intervalo ( 20 ; 30 ] foi o de maior frequência.
Observa -se, ainda, que uma proporção de estudantes menor do que na questão anterior deixou essa questão em branco.
O coeficiente de assimetria é positivo ( 0,48 ), indicando concentração à esquerda e cauda longa à direita.
O mesmo ocorre para \textbf{cada} uma das Grandes Regiões, sendo que a maior assimetria positiva ocorre para a região Centro-Oeste ( com coeficiente 1,54 ). 3.3.2.4 Comentários sobre a correção das respostas à Questão Discursiva 4 A Banca considerou a questão 4 de nível fácil.
Busca o conhecimento do estudante quanto a um assunto básico e fundamental da manutenção industrial, que são os tipos de manutenção aplicados nas indústrias.
Foi considerado que \textbf{está} dentro do contexto de avaliação para o nível de Tecnólogo de Manutenção Industrial.
Quanto ao desempenho dos concluintes, um grande quantitativo não respondeu aos subitens “ d ” e “ e ” ( manutenção baseado na produtividade total e centrada na confiabilidade ), e dos que responderam, poucos mostraram conhecer o assunto.
Uma parcela significativa de respostas não relacionou custos e eficácia, limitando -se apenas a descrever o conceito.
Sendo uma questão com assunto básico, a quantidade de respostas em branco pode ser considerada alta.
Em relação ao tipo de resposta encontrado com maior frequência, destacaram -se respostas que apresentaram somente com os três primeiros tipos de manutenção ; aquelas que não qualificaram os custos em comparação com os demais ; e as que não qualificaram a eficácia comparada com as demais.
As respostas elaboradas pelos estudantes e a própria estrutura da questão, quanto ao que foi solicitado, demonstraram que o espaço destinado para a resposta foi insuficiente.
Mesmo nas respostas com o entendimento correto dos tipos de manutenção, os discentes mostraram dificuldade em escrever.
A forma de descrição, se levada ao pé da letra, implicaria uma redução da média em todas as provas respondidas. 3.3.2.5 Análise da Questão Discursiva 5 do Componente de Conhecimento Específico A Tabela 3.14 contém as informações relativas à questão 5 do conjunto do componente de Conhecimento Específico.
A nota média dos estudantes de todo o Brasil foi 21,2.
A maior média foi registrada na região Nordeste ( 28,8 ), enquanto a menor média foi registrada na região Norte ( 18,4 ).
Quanto à variabilidade das notas, o desvio padrão dos alunos do Brasil, como um todo, foi 21,5.
O maior desvio foi encontrado na região Nordeste ( 24,5 ), e o menor, na região Centro-Oeste ( 18,7 ).
A nota máxima ( 80,0 ) foi alcançada nas regiões Norte, Nordeste e Sudeste.
Nas regiões Sul e Centro-Oeste, a nota máxima foi 70,0.
A nota mínima foi zero para todas as regiões do Brasil.
A menor mediana foi a da região Norte ( 0,0 ) e a maior ( 30,0 ) ocorreu nas regiões Nordeste e Centro-Oeste.
O Gráfico 3.21 representa a distribuição das notas da questão discursiva 5 no componente de Conhecimento Específico.
Destaca -se o grande número de estudantes que deixaram esta questão em branco ou obtiveram nota zero tendo respondido a questão.
Nota -se, ainda, concentração de notas no intervalo ( 20 ; 30 ] e poucos alunos com notas superiores a este intervalo.
Como nas demais questões discursivas do componente de Formação Especifica, a distribuição de notas do Brasil, assim como a distribuição de \textbf{cada} uma das regiões, é positivo ( 0,75 ).
O maior coeficiente de assimetria é o da região Norte ( 1,08 ), e o da região Centro-Oeste, o menor ( 0,16 ). 3.3.2.6 Comentários sobre a correção das respostas à Questão Discursiva 5 Trata -se de questão também fácil, que busca o conhecimento do estudante quanto a um assunto básico e fundamental da manutenção, os elementos de fixação, que são encontrados em todos os tipos de máquinas.
A banca de corretores considerou que \textbf{está} dentro do contexto de avaliação para o nível de Tecnólogo de Manutenção Industrial.
No entanto, uma boa quantidade dos discentes não entendeu o que \textbf{estava} sendo pedido e não fez a relação entre inovação e quantidade menor de parafusos no redutor mais moderno.
Somente um respondente mencionou diminuição das peças sobressalentes em estoque.
A grande maioria não respondeu ao solicitado na questão de uma forma clara.
Muitos somente sugeriram adotar o mesmo padrão de número de parafusos para ambos os redutores.
O desempenho dos discentes, portanto, foi fraco.
O tipo de resposta encontrado com maior frequência foi a comparação de diferenças entre os dois tipos de redutores, porém sem relacionar com a inovação.
Também foram recorrentes respostas sem qualificar a diminuição do custo da manutenção e a diminuição de sobressalentes em estoque. 3.3.3 Considerações Finais O trabalho de correção das provas do ENADE/2011 dos alunos de Tecnologia em Manutenção Industrial foi bastante facilitado pela fixação, tanto dos gabaritos e conteúdos, quanto das pontuações norteadoras, em acordo com a Comissão Assessora do INEP.
O fato de realizar a correção por meio do Sistema Online favoreceu sobremaneira a tarefa, com destaque para a eficiência não somente do Sistema em si, mas de todos os profissionais da fundação organizadora.
Em seu conjunto, a prova solicitou dos estudantes conceitos fundamentais, o estabelecimento de relações lógicas e conhecimento geral sobre temas essenciais na Área, diretamente relacionados à atuação profissional.
Entretanto, o desempenho dos estudantes, de maneira geral, ficou abaixo do esperado, demandando reflexões a respeito de mudanças no processo de formação.

% matched lemmas: cada, certo, estar, ordem
\end{textsample}
